\chapter{Bowel Incontinence}

\section*{Definition}
Fecal incontinence is involuntary leakage of stool or mucus, ranging from occasional staining to complete loss of bowel control. It can affect skin health, confidence, and quality of life.

\section*{When to See a Doctor}

\begin{itemize}
  \item Arrange assessment for new or persistent leakage, especially with diarrhea, constipation, rectal bleeding, weight loss, or reduced quality of life.
  \item Seek emergency care for sudden loss of bowel control with new leg weakness, numbness around the saddle area, or bladder dysfunction.
\end{itemize}

\section*{How It Develops}
Continence depends on the anal sphincters, rectal sensation and storage, stool consistency, and pelvic-floor and nerve function. Leakage may result from urgency, passive soiling, overflow from impaction, diarrhea, sphincter injury, or impaired sensation.

\section*{Causes}
\begin{itemize}
 \item Obstetric or surgical injury to the anal sphincter or rectum
 \item Neurological conditions, including stroke, multiple sclerosis, Parkinson disease, or spinal nerve disorders
 \item Chronic constipation with rectal impaction and overflow leakage
 \item Pelvic-floor muscle weakness or impaired rectal sensation
 \item Diarrhea, inflammatory bowel disease, infection, or medicines that loosen stool
\end{itemize}

\section*{Related Symptoms}
\begin{itemize}
 \item Urgency, constipation or diarrhea, abdominal discomfort, skin irritation, and perianal rash
\end{itemize}


\section*{Medical Evaluation}
\begin{itemize}
 \item History assesses stool consistency, urgency, leakage pattern, obstetric and surgical history, neurologic symptoms, medicines, and constipation or diarrhea.
 \item Examination may include the abdomen, perianal skin, digital rectal examination, and assessment of sphincter and pelvic-floor function.
 \item Anoscopy, endoscopy, anorectal manometry, endoanal ultrasound, or pelvic-floor testing is selected when the history and examination indicate it.
\end{itemize}

\section*{Treatment Options}
\begin{itemize}
 \item Treatment may include fibre or stool regulation, antidiarrheal therapy when appropriate, pelvic-floor exercises and biofeedback, skin care, and management of the underlying disease.
 \item Persistent or severe incontinence may require specialist treatment, such as biofeedback, sacral nerve stimulation, repair of a sphincter defect, or selected surgery. Injectable bulking agents have limited and variable benefit.
\end{itemize}

\section*{Self-Care and Home Management}
\begin{itemize}
 \item Keep a bowel diary, including stool consistency, urgency, leakage, diet, and medicines.
 \item Keep the perianal skin clean and dry; use gentle cleansing and a barrier cream if irritation develops.
 \item Regular toileting, adequate fluids, and individualized soluble fibre may help, but fibre should be adjusted if constipation or impaction is present.
 \item Do not self-treat persistent leakage with loperamide without medical advice, especially if there is constipation, blood, fever, suspected infection, or inflammatory bowel disease.
\end{itemize}

\section*{References}
\begin{thebibliography}{99}
\bibitem{bowel_incontinence:ref1} Bharucha AE, Knowles CH, et al. ``Fecal incontinence in adults.'' \textit{Nat Rev Dis Primers}. 2022;8(1):53.
\bibitem{bowel_incontinence:ref2} Whitehead WE, Wald A, et al. ``Fecal incontinence.'' \textit{Gastroenterology}. 2004;126(1 Suppl 1):S40-S45.
\bibitem{bowel_incontinence:ref3} American College of Gastroenterology. ``ACG clinical guideline: management of fecal incontinence.'' \textit{Am J Gastroenterol}. 2022;117(9):1400-1416.
\bibitem{bowel_incontinence:ref4} Feldman M, Friedman LS, Brandt LJ, eds. \textit{Sleisenger and Fordtran\'s Gastrointestinal and Liver Disease}, 11th ed. Elsevier, 2021.
\bibitem{bowel_incontinence:ref5} National Institute of Diabetes and Digestive and Kidney Diseases. ``Definition and Facts of Fecal Incontinence.'' Accessed September 2026. https://www.niddk.nih.gov/health-information/digestive-diseases/bowel-control-problems-fecal-incontinence/definition-facts.
\bibitem{bowel_incontinence:ref6} American Society of Colon and Rectal Surgeons. ``Clinical Practice Guidelines for the Treatment of Fecal Incontinence.'' 2023. https://fascrs.org/ascrs/media/files/2023-Fecal-Incontinence-CPG.pdf.
\end{thebibliography}
